\section{Microphone Arrays: Spatial Methods for DOA Estimation}

\subsection{Introduction}
In this chapter, we focus on spatial methods for \textbf{Direction of Arrival (DOA) estimation} using \textbf{microphone arrays}.
The goal is to detect and locate $N$ radiating \textbf{acoustic sources} in space by leveraging an array of $M$ passive microphones.
Rather than directly identifying the sources, we estimate the acoustic energy distribution across space by means of \textbf{spatial spectral analysis}.
This approach provides information about the directions from which the sound arrives, allowing us to infer the positions of the sources.

We assume a \textbf{simplified scenario} where the sources are located in the \textbf{far-field}, such that wavefronts can be approximated as planar, so that each source can be characterized solely by its direction of arrival (DOA).
Additionally, we assume a \textbf{2D geometry} in which all sources and sensors lie on the same plane, neglecting reverberation, dispersion and any propagation inhomogeneities.

\begin{figure}[H]
    \centering
    \includegraphics[width=0.5\linewidth]{DOA scenario.png}
    \caption{DOA estimation in a simplified scenario}
\end{figure}

\subsection{Data model}
For the moment, let us consider a single source radiating from a certain DOA $\theta$.
The output of the $k$-th sensor can be modelled as:

\[
y_k(n) = h_k(n) \ast s(n - \tau_k) + e_k(n), \qquad k = 1, 2, \cdots, M
\]

The signal $s$ is the \textbf{source signal}, as measured at a reference point, while $h_k$ is the \textbf{impulse response} of the $k$th sensor.
The time $\tau_k$ is the \textbf{propagation delay} between the reference point and the $k$-th sensor, while the signal $e_k$ is the \textbf{additive noise} at the sensor.

The information about the source position $\theta$ is carried by the \textbf{propagation delays} $\tau_k$.
The basic problem is to \textbf{estimate} these time delays, assuming the impulse responses of the sensors are known.

The model equation significantly simplifies if we consider an \textbf{emitted narrowband signal}, with known carrier frequency $\omega_c$.
In this case, the signal $s$ can be modelled as:

\[
s(n) = \alpha(n) \cos \left[\omega_c n + \phi(n) \right]
\]

Note: Anna and Peo remember you that we are operating in discrete domain, so $\omega_c$ is a normalized frequency.

We extract frames $s_t$ of length $L$ samples from the signal using rectangular window $w$ and unitary hop-size:

\[
s_t(n) = s(n) \ w(n - t), \qquad
w(n) = 
\begin{cases}
1, & 0 \leq n \leq L - 1 \\
0, & \text{otherwise}
\end{cases}
\]

\clearpage

For each frame, we compute the Discrete-Time Fourier Transform (DTFT):

\[
S_t(\omega) = \sum_{n=-\infty}^{\infty} s_t(n) \ e^{-j \omega n}
\]

Notice that this corresponds to the STFT of $s(n)$, computed with unitary time resolution.

Considering the DTFT at the source frequency $\omega = \omega_c$:

\[
S_t(\omega_c)
= \sum_{n=-\infty}^{\infty} s_t(n) \ e^{-j \omega_c n}
= \sum_{n=-\infty}^{\infty} s(n) \ w(n - t) \ e^{-j \omega_c n}
= \sum_{n=t}^{t+L-1} s(n) \ e^{-j \omega_c n}
= \sum_{n=t}^{t+L-1} \alpha(n) \cos\left[\omega_c n + \phi(n)\right] e^{-j \omega_c n}
\]

\[
\begin{array}{cc}
    n = l + t\\
    l = n - t 
\end{array}
\longrightarrow \quad
= \sum_{l=0}^{L-1} \alpha(l + t) \cos\left[\omega_c(l + t) + \phi(l + t)\right] e^{-j \omega_c(l + t)}
= e^{-j \omega_c t} \sum_{l=0}^{L-1} \alpha(l + t) \cos\left[\omega_c(l + t) + \phi(l + t)\right] e^{-j \omega_c l}
\]

We assume that the amplitude $\alpha(n)$ and the phase $\phi(n)$ are \textbf{slowly-varying} functions over time, which means that over a short observation window of length $L$, they can be considered approximately constant:

\[
\alpha(t) \approx \alpha(t + l), \qquad \phi(t) \approx \phi(t + l), \qquad \forall \ l = 0, 1, \cdots, L - 1
\]

The STFT expression can be simplified accordingly:

\[
S_t(\omega_c)
= e^{-j \omega_c t} \sum_{l = 0}^{L - 1} \alpha(t) \cos\left[\omega_c(l + t) + \phi(t)\right] e^{-j \omega_c l}
\]

\[
= e^{-j \omega_c t} \ \alpha(t) \sum_{l = 0}^{L - 1} \cos\left[\omega_c l + \omega_c t + \phi(t)\right] e^{-j \omega_c l}
\approx e^{-j \omega_c t} \ \alpha(t) e^{j\left(\omega_c t + \phi(t)\right)}
= \alpha(t) \ e^{j \phi(t)}
\]

As the signal is narrowband and centered at the known frequency $\omega_c$, the previous expression can be interpreted as a function of the time variable only.
To simplify the \textbf{notation}, we redefine the above expression as follows:

\[
s(t) \triangleq S_t(\omega_c) = \alpha(t) e^{j \phi(t)}
\]

From now on, we will consider $s(t)$ as the source signal: we are implicitly working in the STFT domain at a single frequency value corresponding to $\omega = \omega_c$.

The computation of the STFT can also be interpreted as a \textbf{demodulation operation}, followed by a \textbf{low-pass filtering} of the signal.

\begin{figure}[H]
    \centering
    \includegraphics[width=0.42\linewidth]{Bandpass.png}
    \includegraphics[width=0.42\linewidth]{Bandpass demodulation.png}
    \includegraphics[width=0.28\linewidth]{Baseband.png}
    \caption{Band-pass signal (left), demodulated signal (right), base-band signal (bottom)}
\end{figure}

We start from the band-pass (modulated) signal:

\[
\overline{s}(n) = \alpha(n) \cos\left[ \omega_c n + \phi(n) \right]
= \frac{1}{2} \alpha(n) \left[ e^{j \left( \omega_c n + \phi(n) \right)} + e^{-j \left( \omega_c n + \phi(n) \right)} \right]
\]

Then, we demodulate by multiplying with $2e^{-j\omega_c n}$:

\[
\tilde{s}(n) = \overline{s}(n) \cdot 2e^{-j\omega_c n}
= \alpha(n) \left[ e^{j\phi(n)} + e^{-j(2\omega_c n + \phi(n))} \right]
\]

Finally, we apply a low-pass filter to remove the high-frequency component and recover the base-band signal:

\[
s(n) = \alpha(n) \ e^{j\phi(n)}
\]

\subsection{Array model}
Using the same \textbf{notation} as before, the STFT of the microphone signals can be expressed as:

\[
y_k(t) = H_k(\omega_c) \cdot \underbrace{s(t) \ e^{-j\omega_c \tau_k}}_{\text{STFT of } s(n-\tau_k)} + e_k(t)
\]

The term $H_k(\omega_c)$ is the frequency response of the $k$-th microphone evaluated at the frequency $\omega = \omega_c$.
To make the model more compact, we adopt a vector representation and derive the complete \textbf{array model for a single source}.
The microphone array output becomes:

\[
\underline{\mathrm{y}}(t) = \underline{\mathrm{a}}(\theta) \ s(t) + \underline{\mathrm{e}}(t), \qquad
\underline{\mathrm{y}}(t) =
\begin{bmatrix}
    y_1(t) \\
    y_2(t) \\
    \vdots \\
    y_M(t)
\end{bmatrix}
,\quad
\underline{\mathrm{a}}(\theta) =
\begin{bmatrix}
    H_1(\omega_c) \ e^{-j\omega_c \tau_1} \\
    H_2(\omega_c) \ e^{-j\omega_c \tau_2} \\
    \vdots \\
    H_M(\omega_c) \ e^{-j\omega_c \tau_M}
\end{bmatrix}
,\quad
\underline{\mathrm{e}}(t) =
\begin{bmatrix}
    e_1(t) \\
    e_2(t) \\
    \vdots \\
    e_M(t)
\end{bmatrix}
\]

We call $\underline{\mathrm{y}}$ the \textbf{array vector}, $\underline{\mathrm{a}}$ the \textbf{propagation vector} and $\underline{\mathrm{e}}$ the \textbf{noise vector}.
We can extend the vector representation to $N$ distinct sources.
Let us define the source signals $s_k$, the corresponding DOAs $\theta_k$ and the propagation vectors $\underline{\mathrm{a}}(\theta_k)$ as:

\[
s_1(t),\ s_2(t),\ \cdots,\ s_N(t)
\qquad \qquad
\theta_1,\ \theta_2,\ \cdots,\ \theta_N
\qquad \qquad
\underline{\mathrm{a}}(\theta_1),\ \underline{\mathrm{a}}(\theta_2),\ \cdots,\ \underline{\mathrm{a}}(\theta_N)
\]

Under the assumption of \textbf{ideal linear acoustic propagation}, we can apply the \textbf{superposition principle} to obtain the array model for the case of multiple sources:

\[
\underline{\mathrm{y}}(t) = \left[\mathrm{A}\right] \underline{\mathrm{s}}(t) + \underline{\mathrm{e}}(t),
\]

We define the following:

\[
\left[\mathrm{A}\right] = 
\begin{bmatrix}
\underline{\mathrm{a}}(\theta_1) & \underline{\mathrm{a}}(\theta_2) & \cdots & \underline{\mathrm{a}}(\theta_N)
\end{bmatrix}_{M \times N}: \text{propagation matrix},
\qquad
\underline{\mathrm{s}}(t) =
\begin{bmatrix}
s_1(t) \\
s_2(t) \\
\vdots \\
s_N(t)
\end{bmatrix}_{N \times 1}: \text{source vector}
\]

\clearpage

\subsection{Uniform Line Array (ULA)}
We now introduce a specific and widely used array geometry: the \textbf{Uniform Linear Array (ULA)}.
This configuration consists of $M$ identical \textbf{omnidirectional sensors} arranged along a \textbf{straight line}, with \textbf{equal spacing} $d$ between consecutive sensors.
The uniform spacing and simple geometry make this setup particularly convenient for DOA estimation tasks.

To simplify the analysis, we consider again the case of a \textbf{single acoustic source} impinging on the array: the signal from the source reaches each microphone with a \textbf{relative delay} that depends on the \textbf{angle of arrival} $\theta$ and the position of each microphone in the array.
Because the wavefront is assumed to be planar (\textbf{far-field assumption}), the time delay between adjacent microphones is constant:

\[
\tau = \frac{d \sin(\theta)}{c}, \qquad c\approx 340~\text{m/s}: \quad \text{speed of sound}
\]

This geometrical configuration allows the propagation model to be expressed in a structured way and forms the basis for many spatial signal processing methods.

\begin{figure}[H]
    \centering
    \includegraphics[width=0.5\linewidth]{ULA.png}
    \caption{Uniform Line Array}
\end{figure}

To proceed with the model derivation, we take the \textbf{first microphone} in the array as the reference point (so without delay).
Under this assumption, we can explicitly express the time delay for the $k$-th microphone as a function of the direction of arrival $\theta$:

\[
\tau_k = (k - 1) \frac{d \sin(\theta)}{c}, \qquad \theta \in [-90^\circ, 90^\circ], \qquad k = 1, 2, \cdots, M
\]


We assume that all microphones in the array are identical and omnidirectional.
Consequently, we can consider all \textbf{microphones frequency responses} to be unitary and equal:

\[
H_1(\omega_c) = H_2(\omega_c) = \cdots = H_M(\omega_c) = 1
\]

By inserting this into the general expression of the propagation vector, we obtain the simplified form:

\[
\underline{\mathrm{a}}(\theta) =
\begin{bmatrix}
    H_1(\omega_c) \ e^{-j\omega_c \tau_1} \\
    H_2(\omega_c) \ e^{-j\omega_c \tau_2} \\
    \vdots \\
    H_M(\omega_c) \ e^{-j\omega_c \tau_M}
\end{bmatrix}
\quad \Rightarrow \quad
\underline{\mathrm{a}}(\theta) = 
\begin{bmatrix}
1 \\
e^{-j \omega_c \frac{d \sin(\theta)}{c}} \\
\vdots \\
e^{-j \omega_c (M-1) \frac{d \sin(\theta)}{c}}
\end{bmatrix}
\]

\subsubsection{Spatial frequency}
The terms of the propagation vector correspond to the samples of a complex sinusoid (that is a plane wave):

\[
e^{-j \omega_c \frac{d \sin(\theta)}{c} k}, \qquad \quad k = 0, 1, \cdots, M-1
\]

The frequency of this sinusoid is called \textbf{spatial frequency}:

\[
\omega_s = \omega_c \frac{d \sin(\theta)}{c}
\]

For a given time instant, the microphones are \textbf{sampling a sinusoid in the spatial domain} (i.e. on a line), whose frequency is proportional to \(\omega_s\) and depends on the source DOA $\theta$.

\begin{figure}[H]
    \centering
    \includegraphics[width=0.75\linewidth]{Spatial frequency.png}
    \caption{Uniform Line Array: spatial sampling for different DOA $\theta$}
\end{figure}

\subsubsection{Anti-aliasing condition}
The complex sinusoid in the propagation vector $\underline{a}(\theta)$ is sampled with \textbf{unitary sampling frequency} $k$:

\[
\underline{a}(\theta) = \left[ 1 \quad e^{-j\omega_s} \quad \cdots \quad e^{-j(M-1) \ \omega_s} \right]^T
\]

For the sampling theorem to hold, we must satisfy:

\[
|\omega_s| \leq \pi
\]

The anti-aliasing condition can be expressed as a function of the microphone spacing $d$:

\[
\omega_s = \omega_c \frac{d \sin(\theta)}{c} = 2 \pi \frac{f_c}{c} \cdot d \sin(\theta) = 2\pi \frac{d \sin(\theta)}{\lambda}
\quad \Rightarrow \quad
\abs{ 2\pi \frac{d \sin(\theta)}{\lambda}} \leq \pi
\]

\[
\Rightarrow \quad d \abs{\sin(\theta)} \leq \frac{\lambda}{2} \quad \Rightarrow \quad d \leq \frac{\lambda}{2}
\]

\subsection{DOA estimation problem}
The \textbf{source DOA} is directly related to the \textbf{spatial frequency}.
Thus, DOA estimation can be brought to the problem of spatial frequency estimation: $N$ sources impinging (interfering) from different DOAs are characterized by different spatial frequencies.

\textbf{Non-parametric methods} (i.e. spatial filtering) perform DOA estimation without any assumptions about the covariance structure of the array data.
\textbf{Parametric methods} perform DOA estimation based on a parametrization of the \textbf{covariance structure} of the array data.

\subsection{Non-parametric methods: Spatial filtering}
The idea is to apply a \textbf{linear combination} of the microphone signals in order to enhance the signal coming from a given direction $\theta$ and to attenuate the signal coming from undesired directions.
Using the weights $\underline{\mathrm{h}} = \left( h_{1}, h_{2}, \dots, h_{M} \right)^{H}$, the \textbf{spatial-filtered signal} is obtained as:

\[
y_{F}(t) = \sum_{k=1}^{M} h_{k} \cdot y_{k}(t) = \underline{\mathrm{h}}^{H} \ \underline{\mathrm{y}}(t)
\]

Assuming that the \textbf{filter} $\underline{\mathrm{h}}$ is \textbf{band-pass} and centered at $\theta$ in the spatial domain, the power of the filtered signal should give a good indication of the energy coming from the direction $\theta$.
The \textbf{power} of the filtered signal is given by:

\[
E\left\{ \abs{y_{F}(t)}^{2} \right\} = \underline{\mathrm{h}}^{H} \left[\mathrm{R}\right] \ \underline{\mathrm{h}},
\qquad
\left[\mathrm{R}\right] = E\left\{ \underline{\mathrm{y}}(t) \underline{\mathrm{y}}^{H}(t) \right\}: \text{covariance matrix of array data}
\]

\subsubsection{Analogy with temporal filtering}
Let us make an analogy between temporal filtering and spatial filtering.
In the temporal filtering of a sinusoidal input at frequency $\omega$, the temporal samples $u(t), u(t-1), \cdots, u(t-M+1)$ are multiplied by filter coefficients $h_{0}, h_{1}, \cdots, h_{M-1}$ and summed to produce the filtered signal:

\[
y_{F}(t) = \sum_{k=0}^{M-1} h_{k} \cdot u(t-k) = \underbrace{\left[ \underline{\mathrm{h}}^{H} \underline{\mathrm{a}}(\omega) \right]}_{\text{frequency response}} \ u(t), \qquad \underline{\mathrm{a}}(\omega): \text{frequency steering vector}
\]


\begin{figure}[H]
    \centering
    \includegraphics[width=0.5\linewidth]{Temporal sampling.png}
    \includegraphics[width=0.47\linewidth]{Spatial sampling.png}
    \caption{Analogy between temporal sampling (left) and spatial sampling (right)}
\end{figure}

In the spatial filtering of a single narrowband source with DOA equal to $\theta$, the spatial samples $y_{1}(t), y_{2}(t), \cdots, y_{M}(t)$ are multiplied by filter coefficients $h_{1}, h_{2}, \cdots, h_{M}$ and summed to produce the filtered signal:

\[
y_{F}(t) = \sum_{k=1}^{M} h_{k} \cdot y_{k}(t) = \underbrace{\left[ \underline{\mathrm{h}}^{H} \underline{\mathrm{a}}(\theta) \right]}_{\text{spatial response}} \ s(t), \qquad \underline{\mathrm{a}}(\theta): \text{spatial steering vector}
\]

\subsubsection{Delay-and-sum beamformer}
The goal behind the \textbf{delay-and-sum beamformer} is to design a \textbf{spatial band-pass filter} $\underline{\mathrm{h}}(\overline{\theta})$ such that it passes undistorted the signals with a target DOA $\overline{\theta}$ and attenuates all the other DOAs different from $\overline{\theta}$ as much as possible.
The first condition is met when the spatial response at $\overline{\theta}$ is unitary:

\[
\underline{\mathrm{h}}^{H}(\overline{\theta}) \ \underline{\mathrm{a}}(\overline{\theta}) = 1
\]

Assuming that the signal is \textbf{spatially white}, that is $\left[\mathrm{R}\right] = \left[\mathrm{I}_{M}\right]$, minimizing the power from all the directions but $\overline{\theta}$ means solving the following problem:

\[
\underline{\mathrm{h}}(\overline{\theta}) = \arg \min_{\underline{\mathrm{h}}} \ \underline{\mathrm{h}}^{H} \ \underline{\mathrm{h}} \quad \Rightarrow \quad
\underline{\mathrm{h}}(\overline{\theta}) = \frac{\underline{\mathrm{a}}(\overline{\theta})}{\underline{\mathrm{a}}^{H}(\overline{\theta}) \ \underline{\mathrm{a}}(\overline{\theta})} = \frac{\underline{\mathrm{a}}(\overline{\theta})}{M}
\]

The \textbf{power} of the filtered signal becomes:

\[
E\left\{ \abs{y_{F}(t)}^{2} \right\} = \underline{\mathrm{h}}^{H}(\overline{\theta}) \left[\mathrm{R}\right] \ \underline{\mathrm{h}}(\overline{\theta}) = \frac{\underline{\mathrm{a}}^{H}(\overline{\theta}) \left[\mathrm{R}\right] \ \underline{\mathrm{a}}(\overline{\theta})}{M^{2}}
\]

The DOAs of the sources can be estimated by evaluating the following \textbf{pseudo-spectrum} function at all the spatial directions:

\[
p(\theta) = \frac{\underline{\mathrm{a}}^{H}(\theta) \left[\hat{\mathrm{R}}\right] \ \underline{\mathrm{a}}(\theta)}{M^{2}}, 
\qquad \theta \in \left[ -\frac{\pi}{2}, \frac{\pi}{2} \right], \qquad \left[\hat{\mathrm{R}}\right] = \frac{1}{K} \sum_{t=1}^{K} \ \underline{\mathrm{y}}(t) \ \underline{\mathrm{y}}^{H}(t)
\]

The matrix $\left[\hat{\mathrm{R}}\right]$ is the \textbf{sample estimate} of the covariance matrix $\left[\mathrm{R}\right]$.
The DOAs of the $N$ sources are estimated as the $N$ highest peaks of $p(\theta)$.

The name delay-and-sum is due to the fact that the resulting filter weights in $\underline{\mathrm{h}}(\overline{\theta})$ correspond to \textbf{pure delays}.
The effect of the filter is to \textbf{re-phase the microphone signals} according to the propagation delays, producing constructive \textbf{interference} for the direction $\overline{\theta}$ and destructive interference for the other directions.
As the filter does not depend on the array data, the spatial response $\underline{\mathrm{h}}(\overline{\theta}) \ \underline{\mathrm{a}}(\overline{\theta})$ can be easily computed.

\begin{figure}[H]
    \centering
    \includegraphics[width=0.5\linewidth]{Beamformer.png}
    \caption{Delay-and sum beamformer: DOA estimation with $\overline{\theta} = 25^{\circ}$ and $M = 10$ microphones}
\end{figure}

The beamformer \textbf{resolution} depends on the width of the main lobe of the spatial response.
For a uniform linear array with \textbf{aperture} $l = (M - 1) d$, the width of the main lobe is approximately given by $\dfrac{\lambda}{l \abs{\cos \overline{\theta}}}$.
Given two sources with \textbf{angular separation} $\Delta \theta$, located around the frontal direction at $0^{\circ}$, they can be resolved by the pseudo-spectrum function $p(\theta)$ whenever the following condition occurs:

\[
\Delta \theta > \frac{\lambda}{l}
\]

In this case, we expect $p(\theta)$ to exhibit two distinct peaks corresponding to the DOAs $\theta_{1}$ and $\theta_{2}$.
If the condition is not satisfied, a single peak will be visible somewhere in the range $\left[ \theta_{1}, \theta_{2} \right]$, leading to inconsistent DOA estimation.

\subsubsection{Capon method}
The Capon method is formulated by the following minimization problem, introduced before:

\[
\underline{\mathrm{h}}(\overline{\theta}) = \arg \min_{\underline{\mathrm{h}}} \ \underline{\mathrm{h}}^{H} \ \left[\mathrm{R}\right] \ \underline{\mathrm{h}},
\qquad
\underline{\mathrm{h}}^{H}(\overline{\theta}) \ \underline{\mathrm{a}}(\overline{\theta}) = 1
\]

Instead of assuming $\left[\mathrm{R}\right] = \left[\mathrm{I}_{M}\right]$, we aim at minimizing the power of signals that actually impinge on the array from directions different from the target one.
In other words, the Capon method puts effort into attenuating the actual interfering (unwanted) sources, rather than attenuating the energy from all the possible locations.
The solution becomes:

\[
\underline{\mathrm{h}}(\overline{\theta}) = \frac{\left[\mathrm{R}\right]^{-1} \ \underline{\mathrm{a}}(\overline{\theta})}{\underline{\mathrm{a}}^{H}(\overline{\theta}) \left[\mathrm{R}\right]^{-1} \ \underline{\mathrm{a}}(\overline{\theta})}
\]

The corresponding output power becomes:

\[
E\left\{ \abs{y_{F}(t)}^{2} \right\} = \underline{\mathrm{h}}^{H}(\overline{\theta}) \left[\mathrm{R}\right] \ \underline{\mathrm{h}}(\overline{\theta}) = \frac{1}{\underline{\mathrm{a}}^{H}(\overline{\theta}) \left[\mathrm{R}\right]^{-1} \ \underline{\mathrm{a}}(\overline{\theta})}
\]

The DOAs are estimated in correspondence to the $N$ highest peaks of the Capon pseudo-spectrum:

\[
p(\theta) = \frac{1}{\underline{\mathrm{a}}^{H}(\theta) \ \left[\hat{\mathrm{R}}\right]^{-1} \ \underline{\mathrm{a}}(\theta)},
\qquad \theta \in \left[ -\frac{\pi}{2}, \frac{\pi}{2} \right]
\]

The filter weights, as well as the spatial response, depend on the data covariance matrix.
The resolution of the Capon method is not constant but depends on the data and the performance is slightly superior to that of beamforming.

\subsection{Parametric methods}
Parametric methods represent a powerful alternative to spatial filtering for DOA estimation.
They leverage a particular \textbf{model of the covariance structure} of the array data.
The number of sources $N$ is exactly known in advance and it should occur that $N < M$.
The source signals $s_{1}(t), s_{2}(t), \cdots, s_{N}(t)$ are such that their covariance matrix $\left[\mathrm{R}_{s}\right]$ is non-singular (partial correlation is admitted):

\[
\mathrm{rank}\left( \left[\mathrm{R}_{s}\right] \right) = N, 
\qquad \left[\mathrm{R}_{s}\right] = E \left\{ \underline{\mathrm{s}}(t) \ \underline{\mathrm{s}}^{H}(t) \right\}
\]

The sensor noise is \textbf{spatially white}, with \textbf{independent and identically distributed (i.i.d.)} components having identical variance:

\[
E \left\{ \underline{\mathrm{e}}(t) \ \underline{\mathrm{e}}^{H}(t) \right\} = \sigma^{2} \left[\mathrm{I}_{M}\right]
\]

The sensor noise is \textbf{uncorrelated} with the source signals.
All the DOAs are different, thus they lead to different spatial frequencies:

\[
\omega_{s_{1}} \neq \omega_{s_{2}} \neq \dots \neq \omega_{s_{N}}
\]

\subsubsection{Covariance model}
Under these assumptions, the \textbf{input covariance matrix} can be decomposed as the sum of two components:
\[
\left[\mathrm{R}\right] = \left[\mathrm{A}\right] \left[\mathrm{R}_{s}\right] \left[\mathrm{A}\right]^{H} + \sigma^{2} \left[\mathrm{I}_{M}\right]
\]

We remind that $\left[\mathrm{A}\right]$ is the propagation vector:

\[
\left[\mathrm{A}\right] = 
\begin{bmatrix}
\underline{\mathrm{a}}(\theta_1) & \underline{\mathrm{a}}(\theta_2) & \cdots & \underline{\mathrm{a}}(\theta_N)
\end{bmatrix}_{M \times N}
\]

Computing the \textbf{eigenvalue decomposition} and sorting the eigenvalues such that $\lambda_{1} \geq \lambda_{2} \geq \cdots \geq \lambda_{M}$, we can rewrite the covariance matrix as:

\[
\left[\mathrm{R}\right] = \left[ \left[\mathrm{U}\right] \ \left[\mathrm{V}\right] \right]
\underbrace{
\begin{bmatrix}
\lambda_{1} & 0 & & \cdots & & 0 \\
0 & \ddots & & \\
& & \lambda_{N} & & & \vdots \\
\vdots & & & \lambda_{N+1} \\
& & & & \ddots & 0 \\
0 & & \cdots & & 0 & \lambda_{M}
\end{bmatrix}}_{\left[\mathrm{\Lambda}\right]}
\begin{bmatrix}
\left[\mathrm{U}\right]^{H} \\
\left[\mathrm{V}\right]^{H}
\end{bmatrix}
\]

The matrix $\left[\mathrm{U}\right]$ contains the orthonormal eigenvectors associated to $\{\lambda_{1}, \dots, \lambda_{N}\}$, while the matrix $\left[\mathrm{V}\right]$ contains the orthonormal eigenvectors associated to $\{\lambda_{N+1}, \dots, \lambda_{M}\}$.

As the source signals and the noise are \textbf{uncorrelated}, we can rewrite the diagonal eigenvalue matrix $\left[\mathrm{\Lambda}\right]$ as:

\[
\left[\mathrm{\Lambda}\right] = \underbrace{\left[\mathrm{\Lambda}_{s}\right]}_{\text{eigenvalues of } \left[\mathrm{A}\right] \left[\mathrm{R}_{s}\right] \left[\mathrm{A}\right]^{H}} + \underbrace{\left[\mathrm{\Lambda}_{e}\right]}_{\text{eigenvalues of } \sigma^{2} \left[\mathrm{I}_{M}\right]} = \left[\mathrm{\Lambda}_{s}\right] + \sigma^{2} \left[\mathrm{I}_{M}\right]
\]

As the spatial frequencies $\omega_{s}$ are all different, the corresponding propagation matrix $\left[\mathrm{A}\right]$ is full rank:

\[
\mathrm{rank}\left( \left[\mathrm{A}\right] \right) = \mathrm{rank}\left( \left[ \underline{\mathrm{a}}(\theta_{1}) \ \underline{\mathrm{a}}(\theta_{2}) \ \dots \ \underline{\mathrm{a}}(\theta_{N}) \right] \right) = N
\]

Moreover, since $\mathrm{rank}\left( \left[\mathrm{R}_{s}\right] \right) = N$, it must be that:

\[
\mathrm{rank}\left( \left[\mathrm{A}\right] \ \left[\mathrm{R}_{s}\right] \ \left[\mathrm{A}\right]^{H} \right) = N
\]

It follows that $\left[\mathrm{A}\right] \left[\mathrm{R}_{s}\right] \left[\mathrm{A}\right]^{H}$ has $M - N$ null eigenvalues:

\[
\left[\mathrm{\Lambda}_{s}\right] =
\begin{bmatrix}
\overline{\lambda}_{1} & & & \\
& \ddots & & \\
& & \overline{\lambda}_{N} & \\
& & & 0 \\
& & & & \ddots \\
& & & & & 0
\end{bmatrix}
\]

Finally, the eigenvalues of $\left[\mathrm{R}\right]$ can be expressed as:

\[
\lambda_{k} =
\left\{
\begin{array}{lc}
\overline{\lambda}_{k} + \sigma^{2}, & k = 1, \cdots, N  \\
\sigma^{2}, & k = N+1, \cdots, M
\end{array}
\right.
\]

Roughly speaking, the eigenvectors in $\left[\mathrm{U}\right]$ are associated with the source signals, while the eigenvectors in $\left[\mathrm{V}\right]$ are associated with noise.
\textbf{Parametric methods} are derived from the properties of $\left[\mathrm{U}\right]$ and $\left[\mathrm{V}\right]$.
These matrices can be obtained \textbf{if and only if} the number of sources $N$ is known and $N < M$, and the noise power is lower than the source strengths.

Regarding the properties of $\left[\mathrm{V}\right]$, the eigenvalue decomposition implies that:

\[
\left[\mathrm{R}\right] \left[\mathrm{V}\right] =
\left[\mathrm{V}\right]
\begin{bmatrix}
\lambda_{N+1} & & \\
& \ddots & \\
& & \lambda_{M}
\end{bmatrix}
=
\left[\mathrm{V}\right]
\begin{bmatrix}
\sigma^{2} & & \\
& \ddots & \\
& & \sigma^{2}
\end{bmatrix}
= \sigma^{2}\left[\mathrm{V}\right]
\]

\clearpage

Replacing the expression of $\left[\mathrm{R}\right]$ we obtain:

\[
\left[\mathrm{A}\right] \left[\mathrm{R}_{s}\right] \left[\mathrm{A}\right]^{H} \left[\mathrm{V}\right] + \sigma^{2} \left[\mathrm{V}\right] = \sigma^{2} \left[\mathrm{V}\right]
\quad \Rightarrow \quad
\left[\mathrm{A}\right] \left[\mathrm{R}_{s}\right] \left[\mathrm{A}\right]^{H} \left[\mathrm{V}\right] = \left[\mathrm{0}\right]_{M \times (M-N)}
\quad \Rightarrow \quad
\left[\mathrm{A}\right]^{H} \left[\mathrm{V}\right] = \left[\mathrm{0}\right]_{N \times (M-N)}
\]

Regarding the properties of $\left[\mathrm{U}\right]$, the eigenvalue decomposition implies that:

\[
\left[\mathrm{R}\right] \left[\mathrm{U}\right] =
\left[\mathrm{U}\right]
\begin{bmatrix}
\lambda_{1} & & \\
& \ddots & \\
& & \lambda_{N}
\end{bmatrix}
=
\left[\mathrm{U}\right]
\begin{bmatrix}
\overline{\lambda}_{1} & & \\
& \ddots & \\
& & \overline{\lambda}_{N}
\end{bmatrix}
+
\left[\mathrm{U}\right]
\begin{bmatrix}
\sigma^{2} & & \\
& \ddots & \\
& & \sigma^{2}
\end{bmatrix}
= \left[\mathrm{U}\right] \left[\mathrm{\overline{\Lambda}}\right] + \sigma^{2} \left[\mathrm{U}\right]
\]

Replacing the expression of $\left[\mathrm{R}\right]$ we obtain:

\[
\left[\mathrm{A}\right] \left[\mathrm{R}_{s}\right] \left[\mathrm{A}\right]^{H} \left[\mathrm{U}\right] + \sigma^{2} \left[\mathrm{U}\right] = \left[\mathrm{U}\right] \left[\mathrm{\overline{\Lambda}}\right] + \sigma^{2} \left[\mathrm{U}\right]
\quad \Rightarrow \quad
\left[\mathrm{A}\right] \left[\mathrm{R}_{s}\right] \left[\mathrm{A}\right]^{H} \left[\mathrm{U}\right] = \left[\mathrm{U}\right] \left[\mathrm{\overline{\Lambda}}\right]
\]

Since $\left[\mathrm{\overline{\Lambda}}\right]$ is non-singular, it can be inverted, leading to:

\[
\left[\mathrm{U}\right] = \left[\mathrm{A}\right] \ \underbrace{\left[\mathrm{R}_{s}\right] \left[\mathrm{A}\right]^{H} \left[\mathrm{U}\right] \left[\mathrm{\overline{\Lambda}}\right]^{-1}}_{\left[\mathrm{C}\right] \ \text{non-singular matrix}}
\]

\subsubsection{Parametric methods: MUSIC (Multiple Signal Classification)}
MUSIC exploits the fact that:

\[
\left[\mathrm{A}\right]^{H} \left[\mathrm{V}\right] = \left[\mathrm{0}\right]_{N \times (M-N)}
\]

This expression can be rewritten in scalar form as:

\[
\left[\mathrm{A}\right]^{H} \left[\mathrm{V}\right] \left[\mathrm{V}\right]^{H} \left[\mathrm{A}\right] = \left[\mathrm{0}\right]
\]

Its solutions are given by the propagation vectors corresponding to the DOAs:

\[
\underline{\mathrm{a}}^{H}(\theta_{k}) \left[\mathrm{V}\right] \left[\mathrm{V}\right]^{H} \underline{\mathrm{a}}(\theta_{k}) = 0,
\qquad k = 1, \cdots, N
\]

The MUSIC algorithm proceeds as follows:
\begin{enumerate}
    \item Estimate the covariance matrix:
    
    \[
    \left[\hat{\mathrm{R}}\right] = \frac{1}{K} \sum_{t=1}^{K} \ \underline{\mathrm{y}}(t) \ \underline{\mathrm{y}}^{H}(t)
    \]
    
    \item From the eigenvalue decomposition, estimate the eigenvectors $\left[\hat{\mathrm{V}}\right]$ related to noise
    
    \item Estimate the DOAs as the $N$ highest peaks of the pseudo-spectrum function:
    
    \[
    p(\theta) = \frac{1}{\underline{\mathrm{a}}^{H}(\theta) \left[\hat{\mathrm{V}}\right] \left[\hat{\mathrm{V}}\right]^{H} \underline{\mathrm{a}}(\theta)},
    \qquad \theta \in \left[ -\frac{\pi}{2}, \frac{\pi}{2} \right]
    \]
\end{enumerate}

\subsubsection{Parametric methods: ESPRIT (Estimation of Signal Parameters via Rotational Invariance Techniques)}
ESPRIT exploits the fact that:

\[
\left[\mathrm{U}\right] = \left[\mathrm{A}\right] \ \underbrace{\left[\mathrm{R}_{s}\right] \left[\mathrm{A}\right]^{H} \left[\mathrm{U}\right] \left[\mathrm{\overline{\Lambda}}\right]^{-1}}_{\left[\mathrm{C}\right] \ \text{non-singular matrix}}
\]

Given $M$ microphones, we split the array into two identical sub-arrays, each containing $\overline{M} = M - 1$ elements:

\begin{figure}[H]
    \centering
    \includegraphics[width=0.4\linewidth]{ESPRIT subarrays.png}
    \caption{ESPRIT splitting of the microphone array}
\end{figure}

From the array data model:

\[
\underline{\mathrm{y}}(t) = \left[\mathrm{A}\right] \underline{\mathrm{s}}(t) + \underline{\mathrm{e}}(t)
\]

We derive the sub-array models:

\[
\left\{
\begin{array}{cc}
\underline{\mathrm{y}}_{1}(t) = \left[\mathrm{A}_{1}\right] \ \underline{\mathrm{s}}(t) + \underline{\mathrm{e}}_{1}(t),
& \left[\mathrm{A}_{1}\right] = \left[ \left[\mathrm{I}_{\overline{M}}\right] \ \left[\mathrm{0}\right] \right] \left[\mathrm{A}\right] \\
\underline{\mathrm{y}}_{2}(t) = \left[\mathrm{A}_{2}\right] \ \underline{\mathrm{s}}(t) + \underline{\mathrm{e}}_{2}(t),
& \left[\mathrm{A}_{2}\right] = \left[ \left[\mathrm{0}\right] \ \left[\mathrm{I}_{\overline{M}}\right] \right] \left[\mathrm{A}\right]
\end{array}
\right.
\]

We may prove that:

\[
\left[\mathrm{A}_{2}\right] = \left[\mathrm{A}_{1}\right] \left[\mathrm{D}\right],
\qquad
\left[\mathrm{D}\right] =
\begin{bmatrix}
e^{-j\omega_{s_{1}}} & & \\
& \ddots & \\
& & e^{-j\omega_{s_{N}}}
\end{bmatrix}
\]

From the matrix $\left[\mathrm{D}\right]$ we can compute the spatial frequencies, which are directly related to the source DOAs: this estimation is our goal.

Consider the covariance matrices $\left[\mathrm{R}_{1}\right]$ and $\left[\mathrm{R}_{2}\right]$ for the two sub-arrays.
From their eigenvalue decompositions, we can compute the eigenvectors $\left[\mathrm{U}_{1}\right]$ and $\left[\mathrm{U}_{2}\right]$ associated with the sources, which can also be computed as:

\[
\left\{
\begin{array}{c}
\left[\mathrm{U}_{1}\right] = \left[ \left[\mathrm{I}_{\overline{M}}\right] \ \left[\mathrm{0}\right] \right] \left[\mathrm{U}\right] \\
\left[\mathrm{U}_{2}\right] = \left[ \left[\mathrm{0}\right] \ \left[\mathrm{I}_{\overline{M}}\right] \right] \left[\mathrm{U}\right]
\end{array}
\right.
\]

Moreover, we know that $\left[ \mathrm{U} \right] = \left[ \mathrm{A} \right] \left[ \mathrm{C} \right]$, 
where $\left[ \mathrm{C} \right]$ is an invertible matrix.  
It follows that:

\[
\left\{
\begin{array}{l}
\left[ \mathrm{U}_1 \right] 
= \left[ \left[ \mathrm{I}_{\overline{M}} \right] \ \left[ 0 \right] \right] \left[ \mathrm{U} \right] 
= \left[ \left[ \mathrm{I}_{\overline{M}} \right] \ \left[ 0 \right] \right] \left[ \mathrm{A} \right] \left[ \mathrm{C} \right] 
= \left[ \mathrm{A}_1 \right] \left[ \mathrm{C} \right] \\
\left[ \mathrm{U}_2 \right] 
= \left[ \left[ 0 \right] \ \left[ \mathrm{I}_{\overline{M}} \right] \right] \left[ \mathrm{U} \right] 
= \left[ \left[ 0 \right] \ \left[ \mathrm{I}_{\overline{M}} \right] \right] \left[ \mathrm{A} \right] \left[ \mathrm{C} \right] 
= \left[ \mathrm{A}_2 \right] \left[ \mathrm{C} \right]
\end{array}
\right.
\quad \Rightarrow \quad
\left\{
\begin{array}{l}
\left[ \mathrm{A}_1 \right] = \left[ \mathrm{U}_1 \right] \left[ \mathrm{C} \right]^{-1} \\
\left[ \mathrm{A}_2 \right] = \left[ \mathrm{U}_2 \right] \left[ \mathrm{C} \right]^{-1}
\end{array}
\right.
\]

Replacing the expression for $\left[ \mathrm{A}_2 \right]$, we obtain:

\[
\left[ \mathrm{A}_2 \right] = \left[ \mathrm{A}_1 \right] \left[ \mathrm{D} \right]
\quad \Rightarrow \quad
\left[ \mathrm{U}_2 \right] 
= \left[ \mathrm{U}_1 \right] \left[ \mathrm{C} \right]^{-1} \left[ \mathrm{D} \right] \left[ \mathrm{C} \right] 
= \left[ \mathrm{U}_1 \right] \left[ \mathrm{\Phi} \right]
\]

The matrix $\left[ \mathrm{\Phi} \right] = \left[ \mathrm{C} \right]^{-1} \left[ \mathrm{D} \right] \left[ \mathrm{C} \right]$ has the same eigenvalues as $\left[ \mathrm{D} \right]$.
The ESPRIT algorithm proceeds as follows:
\begin{enumerate}
    \item Compute the sample covariance matrix $\hat{[ \mathrm{R} ]}$
    
    \item From the eigenvalue decomposition, estimate $\hat{[ \mathrm{U}_1 ]}$ and $\hat{[ \mathrm{U}_2 ]}$
    
    \item Compute the least-squares estimate of $[ \mathrm{\Phi} ]$:
    
    \[
    \hat{[ \mathrm{\Phi} ]} = \left( \hat{[ \mathrm{U}_1 ]}^{H} \hat{[ \mathrm{U}_1 ]} \right)^{-1} \hat{[ \mathrm{U}_1 ]}^{H} \hat{[ \mathrm{U}_2 ]}
    \]
    
    \item Compute the eigenvalues of $\hat{[ \mathrm{\Phi} ]}$: $\nu_1, \cdots, \nu_N$
    
    \item Estimate the spatial frequencies as:
    
    \[
    \omega_{s_k} = - \arg(\nu_k), \qquad k = 1, \cdots, N
    \]
    
    \item Estimate the DOAs accordingly
\end{enumerate}

\clearpage
